\documentclass[final]{beamer}

% Poster sizing: 78cm x 120cm, portrait
\usepackage[orientation=portrait,size=custom,width=78,height=120,scale=1.4]{beamerposter}
\usepackage{graphicx}
\usepackage{amsmath,amssymb,amsthm}
\usepackage{qrcode}
\usepackage{nicefrac}

% ------------------------------------------------
% Math-like Black & White Theme Configuration
% ------------------------------------------------
\usefonttheme{serif} % Classic LaTeX math font

\setbeamercolor{background canvas}{bg=white}
\setbeamercolor{headline}{bg=white,fg=black}
\setbeamercolor{block title}{bg=white,fg=black}
\setbeamercolor{block body}{bg=white,fg=black}
\setbeamercolor{item}{fg=black}

\setbeamerfont{block title}{series=\bfseries,size=\Large}
\setbeamerfont{headline}{series=\bfseries}

% Headline Template
\setbeamertemplate{headline}{
  \vspace{0.5cm}
  \begin{columns}[c]
    \begin{column}{.2\textwidth}
      \raggedright
      \includegraphics[width=0.9\linewidth]{logos/logo_udd.png}
    \end{column}
    \begin{column}{.6\textwidth}
      \centering
      {\huge \textbf{\inserttitle} \\[1.5ex]}
      {\large \insertauthor \\[1ex]}
      {\normalsize \insertinstitute}
    \end{column}
    \begin{column}{.2\textwidth}
      \raggedleft
      \includegraphics[width=0.9\linewidth]{logos/logo_cics.png}
    \end{column}
  \end{columns}
  \vspace{0.5cm}
  \begin{center}
    \rule{0.95\textwidth}{3pt}
  \end{center}
  \vspace{0.2cm}
}

% Block Template (Minimalist Math Style)
\setbeamertemplate{block begin}{
  \vspace{1ex}
  \begin{center}
    {\usebeamerfont{block title}\MakeUppercase{\insertblocktitle}} \\
    \vspace{0.3ex}
    \rule{0.9\linewidth}{1.5pt}
  \end{center}
  \vspace{1ex}
  \begin{flushleft}
  \large
}
\setbeamertemplate{block end}{
  \end{flushleft}
  \vspace{3ex}
}

% ------------------------------------------------
% Document Info
% ------------------------------------------------
\title{A Cortico-Cerebellar Model for Solving Non-Markovian Environments}
\author{Luis Luarte\inst{1}, \textbf{\textit{Principal Investigator}}\inst{2} (ORCID: 0000-0000-0000-0000)}
\institute{\inst{1} Centro de Investigación en Complejidad Social (CICS), Universidad del Desarrollo, Santiago, Chile \and \inst{2} Department, Center, City, Country}

\begin{document}
\begin{frame}[fragile,t]
  \begin{columns}[T]
    % LEFT COLUMN
    \begin{column}{.45\textwidth}
      
      \begin{block}{Introduction}
        $\vdash$ Cerebellar lesions spare delay conditioning but abolish trace conditioning, revealing a failure to bridge temporal gaps. \\
        \vspace{0.2em}
        $\vdash$ Navigating these non-Markovian gaps requires delayed credit assignment, yet canonical models treat the cerebellum as purely feedforward. \\
        \vspace{0.2em}
        $\vdash$ Feedforward architectures trap the cortex in a Markovian regime restricted to instantaneous observations. \\
        
        \vspace{0.5em}
        \textbf{Hypothesis.} Passing cortical states through a granular expansion, followed by sparse compression and thalamic return, provides the necessary temporal integration to resolve non-Markovian delays without backpropagation-through-time.
      \end{block}

      \begin{block}{Results}
        \vspace{0.5cm}
        \begin{figure}
            \centering
            \includegraphics[width=0.85\linewidth,keepaspectratio]{../cortico-cerebellar-model/presentations/plots/plot_nsga2_hv.png}
            \caption{\textbf{NSGA-II Discovery.} Multi-objective evolution jointly minimized distillation and empirical loss to converge on structural bounds.}
        \end{figure}
        
        \vspace{0.8cm}
        \begin{figure}
            \centering
            \includegraphics[width=0.85\linewidth,keepaspectratio]{../cortico-cerebellar-model/presentations/plots/plot_psis_loo.png}
            \caption{\textbf{Zero-Shot Transfer.} The distilled physiological model predicts raw behavior significantly better than natively optimized heuristics.}
        \end{figure}
      \end{block}
    \end{column}
    
    % RIGHT COLUMN
    \begin{column}{.45\textwidth}

      \begin{block}{Results (cont.)}
        \vspace{0.5cm}

        
        \vspace{0.8cm}

      \end{block}

      \begin{block}{Methods}
        \textbf{State Expansion \& Fading Memory.} The granular layer acts as an Echo State Network to generate an eligibility trace.
        \begin{equation}
          \phi(X_t) = W_{\text{exp}} X_t
        \end{equation}
        Where $X_t$ is the cortical state, $W_{\text{exp}}$ is the fixed immersion matrix, and $\phi(X_t)$ is the expanded trace.
        
        \vspace{0.5em}
        \textbf{Sparse Compression.} The DCN compresses the trace without geometric collapse.
        \begin{equation}
          C_t = P_\theta(\phi(X_t))
        \end{equation}
        Where $P_\theta$ is a sparse Achlioptas projection preventing signal cross-talk, and $C_t$ is the compressed historical state returned via the thalamus.

        \vspace{0.5em}
        \textbf{Teacher-Student Distillation.} The ODE is evolved via NSGA-II.
        \begin{equation}
          \mathcal{J}(m) = \mathbb{E}_{\text{Data}} [\log p(\text{RT}, A \mid m)] - \beta \cdot I(\phi(X_t); C_t)
        \end{equation}
        Where $\mathcal{J}(m)$ optimizes behavioral fit while penalizing mutual information to enforce sparsity.
      \end{block}

      \begin{block}{Conclusion}
        $\vdash$ Fading memory and cortico-thalamic feedback loops are sufficient to resolve non-Markovian environments. \\
        \vspace{0.2em}
        $\vdash$ Teacher-Student distillation via NSGA-II successfully isolates chaotic physiological ODE parameters without breaking auto-differentiation chains. \\
        \vspace{0.2em}
        $\vdash$ This framework successfully captures emergent value and uncertainty signals purely through historical tracking.
      \end{block}
      
      \begin{block}{Acknowledgements \& Bibliography}
        \begin{columns}[T]
          \begin{column}{0.6\textwidth}
            \textbf{Acknowledgements.}\\
            We thank the contributors and funding agencies for their ongoing support of computational neuroscience research at UDD.
          \end{column}
          \begin{column}{0.3\textwidth}
            \centering
            % Bibliography QR Code
            \qrcode[height=5cm]{https://github.com/luisluarte/cortico-cerebellar-model} \\
            \vspace{1ex}
            \normalsize Scan for References
          \end{column}
        \end{columns}
      \end{block}

    \end{column}
  \end{columns}
\end{frame}
\end{document}
